% ============================================================================
% SECTION 1.5: DEMOCRATIZATION THROUGH USE CASES (REVISED v3)
% Added: Zero-friction model addition + Budget control examples
% ============================================================================

\section{Democratization Through Adaptive Routing}
\label{sec:use_cases}

To ground our technical contributions in real-world impact, we illustrate how \ours{} removes two barriers that prevent access to AI capabilities: \emph{cost barriers} (frontier models are prohibitively expensive) and \emph{expertise barriers} (existing routing tools require ML knowledge and extensive setup). Each scenario demonstrates how BanditGPT makes adaptive routing accessible to users who previously lacked either budget or technical infrastructure.

\subsection{Student Projects: Learning Without Prerequisites}

\paragraph{The Dual Barrier.} A student building a paper summarization tool (5,000 abstracts) faces prohibitive costs (\$21.90 for GPT-4o-only) and lacks the hundreds to thousands of labeled examples typically required to calibrate FrugalGPT~\cite{chen2023frugalgpt}.
 Second, implementing FrugalGPT requires hundreds to thousands labeled examples to calibrate chains and expertise to design scoring functions. The student has neither budget nor training data.

\paragraph{With BanditGPT.} The student installs the library via pip, downloads pre-trained priors (1MB), and writes 10 lines of Python:

\begin{small}
\begin{verbatim}
from banditgpt import Router
router = Router(
    mode="standard", 
    lambda_cost=10,
    max_budget_total=50.00  # Semester budget limit
)
results = [router.query(abstract) for abstract in dataset]
\end{verbatim}
\end{small}

The router autonomously identifies that 72\% of abstracts can be handled by Gemini~Flash~Lite (\$0.10/1k), routing only complex technical content to GPT-4o. Total cost: \textbf{\$3.50} (84\% reduction). Setup time: \textbf{5 minutes}. Required expertise: \textbf{None}---no labeled data, no scorer design, no chain configuration.

\paragraph{Budget Control Without Complexity.}
The student specifies their semester budget upfront via \texttt{max\_budget\_total=50.00}. The router tracks cumulative spend and automatically shifts to cheaper models as the budget depletes, ensuring the student doesn't exhaust funds before semester end. No manual model switching; no spreadsheet tracking; no budget overruns.

\paragraph{Adding New Models Without Friction.}
Midway through the semester, Llama-3.3-70B releases at \$0.88/1k---cheaper than GPT-4o (\$4.38/1k) with strong performance. With FrugalGPT, the student would need to:
\begin{enumerate}[leftmargin=*,nosep]
    \item Re-run 500 test queries through Llama-3.3 (\$0.44)
    \item Grade outputs with GPT-4o-as-judge (\$2.19)
    \item Retrain the scorer (hours of work)
    \item Total: \$2.63 + substantial time investment
\end{enumerate}

With BanditGPT, the student types:
\begin{small}
\begin{verbatim}
router.register_model("llama-3.3-70b", cost=0.88)
\end{verbatim}
\end{small}

Within 50 queries (costing \$0.044 in exploration), the router learns Llama-3.3's strengths and begins routing appropriately. The student immediately benefits from the new model without spending time or money on benchmarking.

\paragraph{Accessibility Impact.} By removing both economic and expertise barriers, \ours{} expands access to AI education from elite institutions (which provide API budgets and graduate-level ML courses) to community colleges and online learners who lack both resources. Students learn to build AI-augmented applications through direct experimentation rather than through prerequisite coursework.

\subsection{Independent Researchers: Methods Without Infrastructure}

\paragraph{The Expertise Barrier Amplified.} A social science researcher analyzing 10,000 interview transcripts for thematic patterns faces not only cost barriers (\$438 for GPT-4o-only) but also methodological barriers. Implementing FrugalGPT requires:
\begin{itemize}[leftmargin=*,nosep]
    \item Annotating 1,000+ examples with "correct" thematic codes (weeks of work)
    \item Training a domain-specific scoring function (requires ML expertise)
    \item Selecting and configuring model chains (requires understanding of LLM capabilities)
    \item Re-running calibration if coding schemes evolve (expensive iteration)
\end{itemize}

For researchers without computational backgrounds, these requirements are insurmountable regardless of budget.

\paragraph{With BanditGPT.} The researcher provides a simple prompt template and cost constraint:

\begin{small}
\begin{verbatim}
router = Router(
    mode="hybrid", 
    lambda_cost=5,
    min_quality=0.85  # Research-grade quality floor
)
codes = router.query_batch(transcripts, 
                           template="Identify themes: {text}")
\end{verbatim}
\end{small}

The system starts with pre-trained priors (no user calibration), autonomously discovers that 88\% of transcripts involve straightforward thematic tagging (routed to Nova-Lite), and reserves GPT-4o for nuanced, ambiguous passages. Total cost: \textbf{\$14.20} (68\% reduction). Setup complexity: \textbf{Zero ML expertise required}.

\paragraph{Quality Assurance Through Constraints.}
The researcher specifies \texttt{min\_quality=0.85} to ensure research-grade output quality. The router automatically restricts exploration to models with high confidence of meeting this threshold, providing methodological rigor without requiring the researcher to understand model capabilities or manually curate chains.

\paragraph{Methodological Shift.} This enables \emph{exploratory analysis} previously reserved for computational social science labs with dedicated ML engineers. Researchers iterate on prompt designs, test multiple coding schemes, and refine their approach---all within discretionary budgets and without needing to learn contextual bandit theory or train BERT-based regressors.

\subsection{Startups: Deployment Without Specialists}

\paragraph{The Infrastructure Barrier.} An early-stage legal tech startup (3 engineers, no ML team) wants to add AI-powered contract review. Beyond the \$52k annual inference cost (GPT-4o-only), implementing FrugalGPT requires:
\begin{itemize}[leftmargin=*,nosep]
    \item Hiring an ML engineer to design and train scoring functions (\$150k+ salary)
    \item Collecting and annotating calibration datasets for each contract type
    \item Maintaining a separate ML pipeline for model updates and chain reconfigurations
    \item Monitoring scorer drift as legal language evolves
\end{itemize}

For a pre-revenue startup with 18-month runway, this infrastructure cost (\$225k total: \$150k salary + \$75k operational overhead) exceeds the inference cost savings, making adaptive routing economically infeasible despite technical merit.

\paragraph{With BanditGPT.} The startup deploys without dedicated ML infrastructure:

\begin{small}
\begin{verbatim}
router = Router(
    mode="hybrid", 
    lambda_cost=3,
    max_cost_per_1k=2.00  # Budget ceiling
)
analysis = router.query(contract, task="clause_review")
\end{verbatim}
\end{small}

Standard mode reduces inference costs to \textbf{\$8.4k annually} (84\% reduction). Hybrid mode provides high-assurance routing (\$16k annually) for risk-sensitive legal queries. Crucially, deployment requires \textbf{zero ML expertise}---the founding team implements adaptive routing with their existing engineering skillset, eliminating the need to hire specialists.

\paragraph{Market-Tracking Without ML Team.}
The startup's model registry evolves monthly:
\begin{itemize}[leftmargin=*,nosep]
    \item Month 1: GPT-4o, Claude-3.5, Gemini-Flash
    \item Month 2: +DeepSeek-V3 (\$0.27/1k, excellent for code)
    \item Month 3: +Llama-3.3 (\$0.88/1k, strong reasoning)
    \item Month 4: +Gemini-2.0-Flash (\$0.10/1k, faster)
\end{itemize}

Engineers add each model in 30 seconds:
\begin{small}
\begin{verbatim}
router.register_model("deepseek-v3", cost=0.27)
\end{verbatim}
\end{small}

The router autonomously evaluates the new model on live traffic and shifts routing patterns if it discovers cost savings. The startup continuously benefits from market evolution without dedicating engineering time to recalibration.

\paragraph{Product Viability.} By removing the infrastructure barrier, \ours{} converts AI from "requires dedicated ML team" to "add via standard engineering workflow." This enables small teams to compete with incumbents who have ML departments, democratizing access to AI-augmented product features.

\subsection{Enterprises: Scale Without Experts}

\paragraph{The Organizational Barrier.} A Fortune 500 customer support organization considering LLM-powered automation (10M queries/year) faces organizational friction. While a centralized ML team \emph{could} implement FrugalGPT, the workflow requires:
\begin{itemize}[leftmargin=*,nosep]
    \item ML team collects calibration data from support team (cross-org coordination)
    \item Weeks of iteration to design domain-specific scoring functions
    \item Ongoing maintenance as support categories evolve (dependencies on ML team)
    \item Separate deployments for each product line (no transfer learning)
\end{itemize}

This organizational overhead delays deployment by 6--12 months, during which executives lose patience and the initiative stalls.

\paragraph{With BanditGPT.} The support engineering team deploys directly without ML dependencies:

\begin{small}
\begin{verbatim}
router = Router(
    mode="hybrid", 
    lambda_cost=1,        # Prioritize reliability
    lambda_latency=0.5,
    min_quality=0.95      # Enterprise SLA
)
\end{verbatim}
\end{small}

Standard mode reduces costs from \$43.8M to \textbf{\$7.0M annually} (84\% reduction). Hybrid mode (\$13.4M) provides selective verification for high-stakes queries (billing disputes, account closures). Deployment time: \textbf{2 weeks} (vs.\ 6--12 months for FrugalGPT setup). Required cross-org dependencies: \textbf{None}---support engineers own the entire pipeline.

\paragraph{Dynamic Multi-Service Allocation.}
The enterprise operates multiple services with different operational requirements:

\begin{small}
\begin{verbatim}
# Customer-facing: quality-first
chatbot = Router(min_quality=0.90, lambda_cost=3)

# Internal code review: cost-first
code = Router(min_quality=0.70, lambda_cost=10)

# Docs generation: balanced
docs = Router(min_quality=0.75, lambda_cost=5)
\end{verbatim}
\end{small}

Each service optimizes independently based on its operational constraints, without requiring separate model pools or manual routing logic. The chatbot team can add new models without coordinating with the code review team, enabling autonomous optimization across departments.

\paragraph{Organizational Transformation.} By removing ML team dependencies, \ours{} enables decentralized AI adoption. Support, sales, HR, and legal departments deploy independently, accelerating time-to-value from quarters to weeks.

\subsection{Cross-Cutting Themes: Two Barriers, One Solution}

Three patterns emerge across these use cases:

\paragraph{1. Ease of Use Unlocks Adoption.}
Cost reduction alone is insufficient if setup requires ML expertise. Students lack training data. Researchers lack computational backgrounds. Startups lack ML teams. By providing \emph{immediate usability} through pre-trained priors and autonomous learning, we expand the user base from "ML practitioners" to "anyone who can write Python."

\paragraph{2. Zero-Friction Evolution Enables Experimentation.}
When adding a new model requires days of benchmarking and \$20--50 in profiling costs, users hesitate to track market evolution. With 30-second registration and zero-cost online learning, users continuously benefit from new releases (DeepSeek-V3, Llama-3.3, Gemini-2.0) without operational friction. This real-time adaptability is critical in a fast-moving market where new models release monthly.

\paragraph{3. Intuitive Control Enables Non-Expert Decision-Making.}
Existing systems require users to understand ML concepts (chain thresholds, classifier weights, exploration rates). \ours{} exposes control through domain terms: \texttt{max\_budget\_total=50}, \texttt{min\_quality=0.90}, \texttt{max\_latency\_ms=500}. A startup founder can specify "I have \$500/month for AI" without understanding Thompson Sampling. This interface design democratizes control, not just deployment.

\subsection{Comparison: Setup Requirements}

\begin{table}[t]
\centering
\small
\caption{\textbf{Accessibility Comparison.} We prioritize ease of deployment to expand the user base beyond ML experts.}
\label{tab:accessibility_comparison}
\resizebox{\columnwidth}{!}{%
\begin{tabular}{lccc}
\toprule
\textbf{Requirement} & \textbf{FrugalGPT} & \textbf{RouteLLM} & \textbf{BanditGPT} \\
\midrule
Calibration Data & 500--2k examples & 1k--5k pairs & 0 (shippable priors) \\
Annotated Labels & Yes (ground truth) & Yes (preferences) & No \\
Scorer Training & Yes (BERT finetune) & Yes (classifier) & No \\
Model Chain Design & Manual (expert) & N/A (2 models) & Autonomous \\
Adding New Models & 1--3 days + \$20--50 & 1--2 days + \$20--50 & 30 sec + \$0 \\
Budget Control & Manual tracking & Not supported & Built-in \\
Setup Time & Days & Hours & Minutes \\
ML Expertise & High & Medium & None \\
\midrule
Target User & ML teams & ML practitioners & Anyone \\
\bottomrule
\end{tabular}%
}
\end{table}

Table~\ref{tab:accessibility_comparison} summarizes the operational requirements. By eliminating calibration data, scorer training, manual chain design, \emph{and} benchmarking requirements for new models, \ours{} reduces the barrier from "requires ML team" to "requires pip install." This shift expands adaptive routing from specialized deployments to mainstream adoption.

\paragraph{Design Philosophy.} These scenarios informed \ours{}'s architecture: pre-trained shippable priors eliminate calibration requirements; autonomous online learning removes expertise dependencies; tunable $\lambda$ parameters support diverse operational contexts; zero-benchmark model registration enables real-time market tracking; intuitive constraint interfaces (budget, quality, latency) democratize control; open-source release ensures zero licensing barriers. The technical contributions in \S\ref{sec:method} exist not as algorithmic novelty for its own sake, but as mechanisms for democratizing access to adaptive routing.


